\documentclass[12pt]{article}
\usepackage{amsmath, amssymb, amsthm}
\usepackage{geometry}
\usepackage{setspace}
\usepackage{hyperref}
\geometry{margin=1in}
\onehalfspacing

\title{Proposal for Senior Capstone Project:\\
\textbf{Mathematical Analysis of the Rendering Equation and its Monte Carlo Approximation}}
\author{Luke Jachimiec\\
Math 401 – Senior Capstone Proposal}
\date{\today}

\begin{document}
\maketitle

\begin{abstract}
Lorem ipsum dolor sit amet, consectetur adipiscing elit. Suspendisse at tristique ligula. Pellentesque eu bibendum diam. Suspendisse tortor eros, malesuada a erat in, sollicitudin fringilla felis. Curabitur suscipit tristique ipsum, at ultrices urna dignissim pharetra. Vestibulum malesuada purus nec nibh interdum mollis. Pellentesque sed venenatis elit, sed mollis mi. Nulla facilisi. Cras sit amet turpis ac urna commodo efficitur ut non massa. Aliquam in arcu iaculis, congue ex et, tincidunt ipsum. Sed a lacus non ligula vestibulum bibendum. Sed malesuada rhoncus congue. Praesent condimentum urna non est eleifend venenatis. Nunc ornare mattis lectus, vel convallis nisl. Suspendisse lobortis mi iaculis, blandit odio vel, scelerisque ligula. Cras a ante lacus. Nullam nisi justo, finibus nec erat vitae, consequat eleifend diam.
\end{abstract}

\section{Background and Introduction}
Light transport in physically based rendering is mathematically modeled by the rendering equation introduced by Kajiya (1986). 
In its standard form, the rendering equation expresses outgoing radiance at a point as an integral over the hemisphere of incoming directions:
\[
L_o(x,\omega_o) = 
\int_{\Omega} 
f_r(x,\omega_i,\omega_o)\,
L_i(x,\omega_i)\,
(\omega_i \cdot n)\,
d\omega_i.
\]

Topics to be introduced in this section:
\begin{itemize}
    \item Radiometric quantities: radiance, irradiance, and the bidirectional reflectance distribution function (BRDF).
    \item The hemisphere $\Omega$ as a subset of the unit sphere $S^2$.
    \item The differential solid angle measure $d\omega = \sin\theta\, d\theta\, d\phi$.
    \item Why closed-form solutions are infeasible except in trivial cases.
\end{itemize}

This section situates the problem in mathematical context and motivates Monte Carlo integration as the primary analytical tool.

\section{Project Description}
The goal of the project is to study the rendering equation from a mathematical perspective and construct Monte Carlo estimators for simplified reflectance models.

This project will:
\begin{itemize}
    \item Provide a rigorous definition of every quantity in the rendering equation, including the domain $\Omega$, the BRDF $f_r$, incoming/outgoing directions $\omega_i, \omega_o$, and the geometry term $(\omega_i \cdot n)$.
    \item Express the integral in both surface-based and parametric forms:
    \[
    \int_{\Omega} g(\omega)\, d\omega
    = \int_0^{2\pi}\int_0^{\pi/2} 
    g(\theta,\phi)\,
    \sin\theta\, d\theta\, d\phi.
    \]
    \item Construct Monte Carlo estimators for the integral:
    \[
    \hat{L}_o = \frac{1}{N}
    \sum_{k=1}^{N} 
    \frac{
        f_r(\omega_k)L_i(\omega_k)(\omega_k\cdot n)
    }{
        p(\omega_k)
    }.
    \]
    \item Analyze variance and convergence using the law of large numbers and central limit theorem.
    \item Compare uniform hemisphere sampling with cosine-weighted importance sampling.
    \item Validate numerically through simple experiments (e.g., Lambertian surfaces with constant illumination).
\end{itemize}

No large-scale software development is included; computational components serve purely to verify mathematical properties.

\section{Foundations}
This section develops the mathematical theory underlying the project.

\subsection*{Mathematical Topics to Address}
\begin{itemize}
    \item Geometry of the hemisphere as a 2-dimensional manifold embedded in $\mathbb{R}^3$.
    \item Change of variables:
    \[
    d\omega = \sin\theta\, d\theta\, d\phi.
    \]
    \item Properties of BRDFs:
    \[
    f_r: \Omega \times \Omega \to \mathbb{R}_{\ge 0}, \qquad
    \int_{\Omega} f_r(\omega_i,\omega_o)(\omega_i\cdot n)\, d\omega_i \le 1.
    \]
    \item Structure of the rendering equation as:
    \begin{itemize}
        \item an integral operator,
        \item a fixed-point equation,
        \item and a linear lighting transport operator.
    \end{itemize}
    \item Monte Carlo integration:
    \[
    \int_{\Omega} g(\omega)\, d\omega 
    = \mathbb{E}\left[\frac{g(X)}{p(X)}\right].
    \]
    \item Variance analysis:
    \[
    \operatorname{Var}(\hat{I})
    = \frac{1}{N}\operatorname{Var}\!\left(\frac{g(X)}{p(X)}\right).
    \]
    \item Importance sampling:
    \[
    p(\omega)\propto g(\omega)
    \quad\Rightarrow\quad
    \text{variance minimized}.
    \]
    \item Derivation of cosine-weighted sampling for Lambertian surfaces:
    \[
    f_r = \frac{\rho}{\pi},\qquad
    g(\omega)=\rho L_i(\omega)\cos\theta,
    \qquad
    p(\omega)=\frac{\cos\theta}{\pi}.
    \]
\end{itemize}

This section will include derivations, diagrams, and analytic reasoning to demonstrate mathematical correctness.

\section{Implementation Plan and Timeline}

\subsection*{Mathematical Work (Primary Focus)}
\textbf{Week 1–2:}
\begin{itemize}
    \item Study radiometric definitions and formalize the rendering equation.
    \item Define all symbols, domains, and measures rigorously.
\end{itemize}

\textbf{Week 3–4:}
\begin{itemize}
    \item Rewrite the integral in parametric coordinates.
    \item Explore simple analytic cases (e.g., uniform illumination).
\end{itemize}

\textbf{Week 5–6:}
\begin{itemize}
    \item Develop Monte Carlo estimators.
    \item Derive uniform and cosine-weighted PDFs.
\end{itemize}

\textbf{Week 7–8:}
\begin{itemize}
    \item Prove convergence properties using expectation and variance.
    \item Compare sampling strategies mathematically.
\end{itemize}


\textbf{Week 9–10:}
\begin{itemize}
    \item Implement small numerical experiments in Python.
    \item Plot convergence curves.
\end{itemize}

\textbf{Week 11–12:}
\begin{itemize}
    \item Summarize numerical and theoretical findings.
\end{itemize}

\textbf{Week 13–14:}
\begin{itemize}
    \item Prepare final paper and presentation.
\end{itemize}

\section{Conclusion}
Lorem ipsum dolor sit amet, consectetur adipiscing elit. Suspendisse at tristique
ligula. Pellentesque eu bibendum diam. Suspendisse tortor eros, malesuada a erat in,
sollicitudin fringilla felis. Curabitur suscipit tristique ipsum, at ultrices urna dignissim
pharetra. Vestibulum malesuada purus nec nibh interdum mollis. Pellentesque sed
venenatis elit, sed mollis mi. Nulla facilisi. Cras sit amet turpis ac urna commodo
efficitur ut non massa. Aliquam in arcu iaculis, congue ex et, tincidunt ipsum. Sed
a lacus non ligula vestibulum bibendum. Sed malesuada rhoncus congue. Praesent
condimentum urna non est eleifend venenatis. Nunc ornare mattis lectus, vel convallis
nisl. Suspendisse lobortis mi iaculis, blandit odio vel, scelerisque ligula. Cras a ante
lacus. Nullam nisi justo, finibus nec erat vitae, consequat eleifend diam.

\section{Bibliography}

\begin{itemize}
    \item Kajiya, James T. “The Rendering Equation.” ACM SIGGRAPH, 1986.
    \item Veach, Eric. \textit{Robust Monte Carlo Methods for Light Transport}. Stanford University, 1997.
    \item Pharr, Matt, Jakob, Wenzel, and Humphreys, Greg. \textit{Physically Based Rendering}. 3rd ed., 2016.
    \item Dutré, Philip. \textit{Global Illumination Compendium}, 2003.
    \item Shirley, Peter. \textit{Fundamentals of Computer Graphics}. CRC Press.
\end{itemize}

\end{document}
